\chapter{القسمة والمضاعفات}\label{ch:g5:division}

القسمة المطوّلة التي تعلّمناها في \cref{ch:g4:division} تنال
ترقيتين: قواسم من رقمين، وخوارج قسمة تمضي بعد النقطة ---
فالعملية $3 \div 4$ تنال أخيرًا جوابًا، هو $0.75$. ويختم الفصل
\omterm{def:g5:division:multiple}{بالمضاعفات} والقواسم وأولى مختصرات القابلية للقسمة.

\section{القسمة على عدد من رقمين}

\begin{method}[القواسم ذات الرقمين]\label{met:g5:division:twodigits}
الطريقة نفسها كما قبل، مع مهارة إضافية: تخمين عدد المرات
التي يدخل فيها القاسم. فاكتب أولًا جدولًا صغيرًا \omterm{def:g5:division:multiple}{لمضاعفات} القاسم.
وفي $986 \div 23$ (\omterm{def:g5:division:multiple}{مضاعفات} 23: 23، 46، 69، 92، 115،
138، 161، 184، 207):
\begin{enumerate}
  \item $98$ عشرة: يدخل $23 \times 4 = 92$، ولا يدخل
        $23 \times 5 = 115$: فالرقم $4$، \omterm{def:g3:division:remainder}{والباقي} $6$؛
  \item وأنزِل $6$: فيصير $66$: ويدخل $23 \times 2 = 46$: فالرقم
        $2$، \omterm{def:g3:division:remainder}{والباقي} $20$؛
  \item \omterm{def:g3:division:remainder}{فخارج القسمة} $42$، \omterm{def:g3:division:remainder}{والباقي} $20$. والتحقّق:
        $23 \times 42 + 20 = 966 + 20 = 986$.
\end{enumerate}
\end{method}

\section{خوارج القسمة العشرية}

\begin{example}[القسمة التي ترفض التوقف عند الباقي]\label{ex:g5:division:decimal}
اِقسم $3$ بيتزات بين $4$ أطفال: للعملية $3 \div 4$ \omterm{def:g3:division:remainder}{خارج قسمة} $0$
وباقٍ $3$ --- ولا فائدة منه! فواصل القسمة بعد النقطة:
\begin{enumerate}
  \item $3$ وحدات $= 30$ عُشرًا؛ و $30 \div 4$: $7$ أعشار، \omterm{def:g3:division:remainder}{والباقي}
        $2$ من الأعشار؛
  \item و $2$ من الأعشار $= 20$ جزءًا من مئة؛ و $20 \div 4 = 5$ أجزاء
        من مئة، \omterm{def:g3:division:remainder}{والباقي} $0$.
\end{enumerate}
إذن $3 \div 4 = 0.75$: فينال كل طفل $0.75$ بيتزا --- أي ثلاثة
أرباع، كما كانت لغة الكسور تعرف من قبل
(\cref{prop:g5:fractions:decimal}).
\end{example}

\begin{method}[مواصلة القسمة بعد النقطة]\label{met:g5:division:continue}
حين تنفد الآحاد ويبقى باقٍ:
\begin{enumerate}
  \item اُكتب \omterm{def:g4:decimals:point}{النقطة العشرية} في \omterm{def:g3:division:remainder}{خارج القسمة}؛
  \item وألحِق صفرًا \omterm{def:g3:division:remainder}{بالباقي} (فتصير الآحاد أعشارًا، والأعشار
        أجزاءً من مئة\,\dots) وواصل القسمة؛
  \item وقف حين يصير \omterm{def:g3:division:remainder}{الباقي} $0$ --- أو حين تبدأ الأرقام
        تتكرر بلا نهاية، مثل $1 \div 3 = 0.333\dots$: فأعطِ عندئذٍ
        قيمة مقرَّبة.
\end{enumerate}
\end{method}

\begin{example}\label{ex:g5:division:continue}
$27 \div 6$: \omterm{def:g3:division:remainder}{خارج القسمة} $4$، \omterm{def:g3:division:remainder}{والباقي} $3$؛ ثم النقطة؛ و $30$ عُشرًا
$\div\ 6 = 5$: إذن $27 \div 6 = 4.5$. \quad
وفي $22 \div 8$: $2$، \omterm{def:g3:division:remainder}{والباقي} $6$؛ ثم $60 \div 8 = 7$ \omterm{def:g3:division:remainder}{والباقي} $4$؛
ثم $40 \div 8 = 5$: إذن $22 \div 8 = 2.75$.
\end{example}

\section{المضاعفات والقواسم}

\begin{definition}[المضاعف والقاسم]\label{def:g5:division:multiple}
\emph{مضاعفات}\index{مضاعف} العدد $6$ هي جدوله ممتدًّا
إلى ما لا نهاية: $0, 6, 12, 18, 24, \dots$ وحين يكون $24$ مضاعفًا للعدد $6$،
نقول أيضًا إن $6$ \emph{قاسم}\index{قاسم} للعدد $24$،
أو إن $24$ \emph{قابل للقسمة} على $6$: فالقسمة
$24 \div 6$ لا تترك باقيًا.
\end{definition}

\begin{proposition}[مختصرات القابلية للقسمة]\label{prop:g5:division:criteria}
بلا قسمة، يكون العدد قابلًا للقسمة:
\begin{itemize}
  \item على $2$ حين يكون رقمه الأخير $0$ أو $2$ أو $4$ أو $6$ أو $8$
        (وهي \omterm{def:g3:numbers:evenodd}{الأعداد الزوجية})؛
  \item وعلى $5$ حين يكون رقمه الأخير $0$ أو $5$؛
  \item وعلى $10$ حين يكون رقمه الأخير $0$؛
  \item وعلى $3$ حين يكون \emph{مجموع أرقامه} قابلًا للقسمة على
        $3$ (مثلًا $741$: $7 + 4 + 1 = 12$، وهو قابل --- وفعلًا
        $741 = 3 \times 247$).
\end{itemize}
\end{proposition}
\admitted

\begin{omfigure}
\begin{tikzpicture}[scale=0.5]
  \foreach \n in {0,...,29} {
    \pgfmathtruncatemacro{\r}{int(\n/10)}
    \pgfmathtruncatemacro{\c}{mod(\n,10)}
    \pgfmathtruncatemacro{\lbl}{\n+1}
    \pgfmathtruncatemacro{\mtwo}{mod(\lbl,2)}
    \pgfmathtruncatemacro{\mthree}{mod(\lbl,3)}
    \ifnum\mtwo=0
      \fill[omDef!25] (\c,-\r) ++(-0.45,-0.45) rectangle ++(0.9,0.9);
    \fi
    \ifnum\mthree=0
      \draw[omThm, thick] (\c,-\r) circle (0.42);
    \fi
    \node[font=\footnotesize] at (\c,-\r) {\lbl};
  }
  \node[font=\small, below] at (4.5,-2.8)
    {المربعات الزرقاء: مضاعفات $2$؛\ \ والدوائر الحمراء: مضاعفات
    $3$؛\ \ وكلاهما: مضاعفات $6$};
\end{tikzpicture}

{\small الأعداد من $1$ إلى $30$: \omterm{def:g5:division:multiple}{مضاعفات} $2$ \omterm{def:g5:division:multiple}{ومضاعفات} $3$ ترسم
أنماطًا --- والأعداد التي تحمل العلامتين معًا ($6$ و $12$ و $18$
و $24$ و $30$) هي بالضبط \omterm{def:g5:division:multiple}{مضاعفات} $6$.}
\end{omfigure}

\begin{example}\label{ex:g5:division:criteria}
هل $534$ قابل للقسمة على $2$؟ ينتهي بالرقم $4$: نعم. وعلى $5$؟ لا. وعلى $3$؟
$5 + 3 + 4 = 12$: نعم. إذن $534$ قابل للقسمة على $6$ أيضًا (على $2$
\emph{وعلى} $3$) --- والتحقّق: $534 = 6 \times 89$.
\end{example}

\section{تمارين}

\begin{exercise}[$\star$]\label{exo:g5:division:1}
احسب القسمات المطوّلة (بجدول \omterm{def:g5:division:multiple}{المضاعفات} أولًا):
$851 \div 23$؛ \quad $704 \div 32$.
\end{exercise}

\begin{exercise}[$\star$]\label{exo:g5:division:2}
احسب \omterm{def:g3:division:remainder}{خوارج القسمة} العشرية الدقيقة: $9 \div 2$؛ \quad $27 \div 4$؛
\quad $33 \div 5$؛ \quad $21 \div 8$.
\end{exercise}

\begin{exercise}[$\star$]\label{exo:g5:division:3}
يتقاسم أربعة أصدقاء فاتورة مطعم قدرها $54$ بالتساوي. فكم يدفع
كل واحد؟ (واصل بعد النقطة.)
\end{exercise}

\begin{exercise}[$\star$]\label{exo:g5:division:4}
احسب $10 \div 3$، مواصلًا ثلاثة أرقام بعد النقطة. فماذا
تلاحظ؟ وأعطِ القيمة مقرَّبة إلى الجزء من مئة.
\end{exercise}

\begin{exercise}[$\star$]\label{exo:g5:division:5}
اُذكر: \omterm{def:g5:division:multiple}{مضاعفات} $7$ حتى $70$؛ وقواسم $18$؛
وقواسم $24$.
\end{exercise}

\begin{exercise}[$\star$]\label{exo:g5:division:6}
أقابل للقسمة على $2$؟ وعلى $5$؟ وعلى $10$؟ وعلى $3$؟ اِختبر كل مختصر على:
$470$؛ \quad $735$؛ \quad $8\,001$؛ \quad $1\,314$.
\end{exercise}

\begin{exercise}[$\star$]\label{exo:g5:division:7}
جِد كل الأعداد بين $60$ و $80$ القابلة للقسمة على
$3$ \emph{وعلى} $5$. (وأي قابلية واحدة يجمعها
ذلك؟)
\end{exercise}

\begin{exercise}[$\star$]\label{exo:g5:division:8}
يُقطع شريط طوله $7.2$ م إلى $6$ قطع متساوية. فكم يبلغ طول كل
قطعة؟ (اِقسم عددًا عشريًّا على صحيح: اِقسم الأمتار، ثم
الأعشار.)
\end{exercise}

\begin{exercise}[$\star$]\label{exo:g5:division:9}
تُرزم $1\,000$ كرة في أكياس من $24$. فكم كيسًا ممتلئًا،
وكم كرة تبقى؟ (بقاسم من رقمين.)
\end{exercise}

\begin{exercise}[$\star\star$]\label{exo:g5:division:10}
أيستطيع $5$ أصدقاء أن يتقاسموا $7$ ألواح شوكولاتة بالتساوي بلا باقٍ؟
أعطِ نصيب كل واحد عددًا عشريًّا. والسؤال نفسه في حالة $7$ أصدقاء
يتقاسمون $5$ ألواح --- فما الذي يختلف؟
\end{exercise}

\begin{exercise}[$\star\star$]\label{exo:g5:division:11}
مستعملًا مختصر مجموع الأرقام، جِد أصغر رقم $d$ يجعل
$52d$ (عددًا من ثلاثة أرقام ينتهي بالرقم $d$) قابلًا للقسمة على $3$. ثم
جِد كل الأرقام $d$ الممكنة.
\end{exercise}
